\section{Reflection symmetry and the two vertex patterns}
\label{geo:section}

Throughout this section \(\Sigma=\partial E\) is a nonstandard
minimizing boundary in \(T\). In particular it is connected, nonplanar,
real analytic, and of genus at least two. We first prove that the
original region has coordinate reflection symmetry, then use stability
to restrict its possible vertex patterns.

\begin{lemma}\label{geo:reflection}
After translating the coordinate origins, the original region \(E\)
is invariant under reflection in each of the six coordinate tori
\(x_i=0,1\).
\end{lemma}

\begin{proof}
Fix a coordinate. Let \(H_a\) be the half-torus between \(x_i=a\) and
\(x_i=a+1\). The occupied volume in \(H_a\) is continuous in \(a\),
and its value at \(a+1\) is the complementary half of the total occupied
volume. Thus some \(H_a\) bisects the occupied volume.

Reflecting each half separately across its two end tori gives two
competitors of the original volume. Their characteristic-function traces
agree across each reflecting seam, so the reflection adds no seam
perimeter, even if the cutting torus is tangent to the original surface.
The intersection of the nonplanar analytic surface with either cutting
torus has zero surface area: positive area would force an analytic open
planar piece and hence a planar connected surface. The average perimeter
of the two competitors is therefore the original minimum. Each competitor
is minimizing and consequently has smooth analytic boundary.

Each open half has both positive occupied volume and positive vacant
volume, so it contains an open patch of the original boundary.
A reflected competitor shares such a patch. This patch cannot be planar,
and Theorem~\ref{pre:structure} makes the competitor's boundary connected.
Two connected closed embedded analytic surfaces sharing an open patch
coincide everywhere. To see that local coincidence is closed, at a limit
point choose analytic graph charts over their common limiting tangent
plane. Their graph functions agree on a nonempty open set, and hence
on the connected chart by the identity theorem.
Connectedness then propagates equality across the whole surface.

The occupied sides also agree. Their outward normals agree on the
unchanged patch; the relative sign of the normal fields is constant on
the connected common boundary. Equivalently, equality of their oriented
boundary distributions makes the difference of the characteristic
functions constant on the connected torus, and agreement on the retained
half makes that constant zero. This reasoning does not require \(h\ne0\).
The original region, not merely a chosen competitor, is therefore
reflection invariant. Repeat in the other two coordinates. Translations
of their origins preserve the symmetries already obtained.
\end{proof}

We may now work in \(Q=[0,1]^3\). Normalized area and volume in \(T\)
are exactly relative perimeter and ordinary volume in \(Q\). The outward
normal transforms equivariantly under each reflection. In particular
\(N_i=0\) on the face \(x_i=0\) or \(1\): the surface meets a face
orthogonally.

\begin{lemma}[Strict coordinate monotonicity]\label{geo:normals}
After possibly reversing individual coordinates in \(Q\),
\begin{equation}
 N_i>0\qquad\text{on }\Sigma\text{ whenever }0<x_i<1.
 \label{geo:strict-normal}
\end{equation}
Consequently \(E\cap Q\) is a down-set for the coordinatewise order.
\end{lemma}

\begin{proof}
On the entire torus surface, \(u=N_i\) satisfies the Jacobi equation
\[
 Lu=0,\qquad L=\Delta_\Sigma+|\mathrm{II}|^2,
\]
because axial translations preserve mean curvature. It is odd under
the corresponding reflection. If it vanished identically, the complete
axial circle field would be tangent to \(\Sigma\); its flow would preserve
\(\Sigma\). The quotient by this free circle action is a closed connected
curve, making \(\Sigma\) a torus, contrary to its genus.

We spell out the constrained nodal-domain argument. The stability form is
\[
 \mathcal Q(f,f)=\int_\Sigma
       \bigl(|\nabla f|^2-|\mathrm{II}|^2f^2\bigr)\,dS\ge0
       \quad\text{if }\int_\Sigma f\,dS=0.
\]
The restriction \(u_D\) to a nodal domain \(D\), extended by zero,
belongs to \(H^1(\Sigma)\) and satisfies
\(\mathcal Q(u_D,u_D)=0\). This follows by testing \(Lu=0\) with
the restriction, or first truncating the signed values at distance
\(\varepsilon\) from zero and passing to the limit.
Restrictions to distinct nodal domains are orthogonal for \(\mathcal Q\).
If there were at least three nodal domains, a nonzero combination \(f\)
of two restrictions could be chosen with zero mean. It would have
\(\mathcal Q(f,f)=0\). Polarizing the stability inequality against all
mean-zero tests gives \(Lf=c\) weakly for a constant \(c\).
On the unused nodal domain \(f=0\), so \(c=0\).
Local elliptic regularity first bootstraps this weak solution to a
smooth one. The Jacobi coefficients are analytic because \(\Sigma\) is,
so analytic elliptic regularity~\cite[Theorem~1.1]{Blatt} and the
identity theorem force \(f=0\) everywhere,
a contradiction. Thus \(u\) has at most two nodal domains.

It vanishes on the two reflection tori. An additional zero in either
open half would force both signs locally: a one-signed Jacobi solution
cannot have an interior zero by the strong minimum principle and unique
continuation. Reflection would give both signs in the other half,
hence at least four nodal domains. Different signs on different pieces
of a half give the same contradiction. Therefore the sign is strict
and fixed on each half. Reverse coordinates to obtain
\eqref{geo:strict-normal}.

An increasing coordinate segment meets \(\Sigma\), when it meets it
in its interior, transversely and only from the occupied to the vacant
side. It cannot cross twice. Following coordinate segments proves the
down-set assertion, also on the faces by continuity.
\end{proof}

This is the graph observation of Ritor\'e, presented with its stability
proof in Ros \cite[Proposition~2.3]{Ros}. The explicit argument above
will also control the exceptional slicing levels.

\begin{proposition}[Vertex patterns]\label{geo:patterns}
No vertex of \(Q\) lies on \(\Sigma\). Up to taking the complementary
region and reversing all coordinates, and up to coordinate permutations,
the only possible vertex sets for a nonstandard minimizer are
\begin{equation}
 \{000,100,010\},\qquad
 \{000,100,010,001\}.
 \label{geo:two-patterns}
\end{equation}
The four-vertex pattern may be placed on the smaller-volume side.
No such assertion is yet made for the three-vertex pattern.
\end{proposition}

\begin{proof}
At a vertex, equivariance under all three reflections would force all
components of a boundary normal to vanish. Thus each vertex lies strictly
on one side. The occupied vertices form a nonempty proper order ideal,
containing \(000\) and excluding \(111\).
Complementation followed by \(x\mapsto(1,1,1)-x\) preserves the down-set
property and replaces the number of occupied vertices by its complement
to eight. We may therefore assume at most four occupied vertices.

For one occupied vertex, the three opposite faces are strictly vacant.
Indeed a point on such a face in the closure of the region would, by
monotonicity, put its nearest axial vertex in the closure as well.
Compactness gives a gap from those faces. Reflect about the occupied
corner to lift the region into a box centered there. Every radial ray
starts inside and crosses the boundary exactly once, transversely:
at a boundary point \(x\ne0\),
\(\sum_i x_iN_i>0\), with signs interpreted by reflection.
The radial function is smooth by the implicit function theorem.
The boundary is a sphere.

For two occupied vertices, they form an edge, say in the \(x\)-direction.
The entire corresponding axial circle is occupied, with a neighborhood,
and the other two opposite faces are separated from the region by gaps.
In the cover \(S^1_x\times\mathbb R^2\), radial rays in the transverse
plane about this circle cross once and transversely, since
\(yN_y+zN_z>0\) at the boundary. A smooth radial function over
\(S^1\times S^1\) parametrizes the boundary, which has genus one.
Both cases are excluded.

Three occupied vertices necessarily give the first pattern in
\eqref{geo:two-patterns}. Four occupied vertices either give the second
pattern or form an entire face. In the latter case that whole face is
strictly occupied and the opposite face strictly vacant. The strict
normal signs give a graph between them, separated from both end faces.
Reflection produces two disjoint graphs over the transverse torus,
contradicting connectedness of \(\Sigma\).

Finally the four-vertex pattern is carried to itself by complementation
and reversal of all coordinates. It may thus be assigned to the
smaller-volume side. Reducing the vertex count alone would not establish
the analogous assertion for the three-vertex pattern; its occupied-side
choice will follow from the flux identity.
\end{proof}


\begin{figure}[ht]
\centering
\begingroup
\setlength{\unitlength}{1.15pt}
\newcommand{\vertexdot}[2]{%
 \put(#1,#2){\color{white}\circle*{6}}%
 \put(#1,#2){\color{black}\circle{6}}}
\newcommand{\patterncube}[1]{%
\begin{picture}(130,112)
 \color{black!55}
 \put(25,20){\line(1,0){60}}
 \put(25,20){\line(1,1){20}}
 \put(85,20){\line(1,1){20}}
 \put(45,40){\line(1,0){60}}
 \put(25,75){\line(1,0){60}}
 \put(25,75){\line(1,1){20}}
 \put(85,75){\line(1,1){20}}
 \put(45,95){\line(1,0){60}}
 \put(25,20){\line(0,1){55}}
 \put(85,20){\line(0,1){55}}
 \put(45,40){\line(0,1){55}}
 \put(105,40){\line(0,1){55}}
 \color{black}
 \vertexdot{25}{20}\vertexdot{85}{20}
 \vertexdot{45}{40}\vertexdot{105}{40}
 \vertexdot{25}{75}\vertexdot{85}{75}
 \vertexdot{45}{95}\vertexdot{105}{95}
 \put(25,20){\circle*{6}}
 \put(85,20){\circle*{6}}
 \put(45,40){\circle*{6}}
 \put(25,12){\makebox(0,0)[t]{\scriptsize \(000\)}}
 \put(85,12){\makebox(0,0)[t]{\scriptsize \(100\)}}
 \put(37,44){\makebox(0,0)[r]{\scriptsize \(010\)}}
 \ifnum#1=4
  \put(25,75){\circle*{6}}
  \put(17,77){\makebox(0,0)[r]{\scriptsize \(001\)}}
 \fi
\end{picture}}
\begin{minipage}{.43\linewidth}
 \centering\patterncube{3}\\[-2pt]
 \small Three occupied vertices
\end{minipage}
\hfill
\begin{minipage}{.43\linewidth}
 \centering\patterncube{4}\\[-2pt]
 \small Four occupied vertices
\end{minipage}
\endgroup
\caption{The two patterns in Proposition~\ref{geo:patterns}.
Solid vertices are occupied and open vertices are vacant.
Only the cube edges are drawn; the boundary surface is not depicted.}
\label{geo:pattern-figure}
\end{figure}


